\documentclass[a4paper]{article}

\def\npart{III}

\def\ntitle{Toric Geometry}
\def\nlecturer{D.\ Ranganathan}

\def\nterm{Michaelmas}
\def\nyear{2019}

\ifx \nauthor\undefined
  \def\nauthor{Qiangru Kuang}
\else
\fi

\ifx \ntitle\undefined
  \def\ntitle{Template}
\else
\fi

\ifx \nauthoremail\undefined
  \def\nauthoremail{qk206@cam.ac.uk}
\else
\fi

\ifx \ndate\undefined
  \def\ndate{\today}
\else
\fi

\title{\ntitle}
\author{\nauthor}
\date{\ndate}

%\usepackage{microtype}
\usepackage{mathtools}
\usepackage{amsthm}
\usepackage{stmaryrd}%symbols used so far: \mapsfrom
\usepackage{empheq}
\usepackage{amssymb}
\let\mathbbalt\mathbb
\let\pitchforkold\pitchfork
\usepackage{unicode-math}
\let\mathbb\mathbbalt%reset to original \mathbb
\let\pitchfork\pitchforkold

\usepackage{imakeidx}
\makeindex[intoc]

%to address the problem that Latin modern doesn't have unicode support for setminus
%https://tex.stackexchange.com/a/55205/26707
\AtBeginDocument{\renewcommand*{\setminus}{\mathbin{\backslash}}}
\AtBeginDocument{\renewcommand*{\models}{\vDash}}%for \vDash is same size as \vdash but orginal \models is larger
\AtBeginDocument{\let\Re\relax}
\AtBeginDocument{\let\Im\relax}
\AtBeginDocument{\DeclareMathOperator{\Re}{Re}}
\AtBeginDocument{\DeclareMathOperator{\Im}{Im}}
\AtBeginDocument{\let\div\relax}
\AtBeginDocument{\DeclareMathOperator{\div}{div}}

\usepackage{tikz}
\usetikzlibrary{automata,positioning}
\usepackage{pgfplots}
%some preset styles
\pgfplotsset{compat=1.15}
\pgfplotsset{centre/.append style={axis x line=middle, axis y line=middle, xlabel={$x$}, ylabel={$y$}, axis equal}}
\usepackage{tikz-cd}
\usepackage{graphicx}
\usepackage{newunicodechar}

\usepackage{fancyhdr}

\fancypagestyle{mypagestyle}{
    \fancyhf{}
    \lhead{\emph{\nouppercase{\leftmark}}}
    \rhead{}
    \cfoot{\thepage}
}
\pagestyle{mypagestyle}

\usepackage{titlesec}
\newcommand{\sectionbreak}{\clearpage} % clear page after each section
\usepackage[perpage]{footmisc}
\usepackage{blindtext}

%\reallywidehat
%https://tex.stackexchange.com/a/101136/26707
\usepackage{scalerel,stackengine}
\stackMath
\newcommand\reallywidehat[1]{%
\savestack{\tmpbox}{\stretchto{%
  \scaleto{%
    \scalerel*[\widthof{\ensuremath{#1}}]{\kern-.6pt\bigwedge\kern-.6pt}%
    {\rule[-\textheight/2]{1ex}{\textheight}}%WIDTH-LIMITED BIG WEDGE
  }{\textheight}% 
}{0.5ex}}%
\stackon[1pt]{#1}{\tmpbox}%
}

%\usepackage{braket}
\usepackage{thmtools}%restate theorem
\usepackage{hyperref}

% https://en.wikibooks.org/wiki/LaTeX/Hyperlinks
\hypersetup{
    %bookmarks=true,
    unicode=true,
    pdftitle={\ntitle},
    pdfauthor={\nauthor},
    pdfsubject={Mathematics},
    pdfcreator={\nauthor},
    pdfproducer={\nauthor},
    pdfkeywords={math maths \ntitle},
    colorlinks=true,
    linkcolor={red!50!black},
    citecolor={blue!50!black},
    urlcolor={blue!80!black}
}

\usepackage{cleveref}



% TODO: mdframed often gives bad breaks that cause empty lines. Would like to switch to tcolorbox.
% The current workaround is to set innerbottommargin=0pt.

%\usepackage[theorems]{tcolorbox}





\usepackage[framemethod=tikz]{mdframed}
\mdfdefinestyle{leftbar}{
  %nobreak=true, %dirty hack
  linewidth=1.5pt,
  linecolor=gray,
  hidealllines=true,
  leftline=true,
  leftmargin=0pt,
  innerleftmargin=5pt,
  innerrightmargin=10pt,
  innertopmargin=-5pt,
  % innerbottommargin=5pt, % original
  innerbottommargin=0pt, % temporary hack 
}
%\newmdtheoremenv[style=leftbar]{theorem}{Theorem}[section]
%\newmdtheoremenv[style=leftbar]{proposition}[theorem]{proposition}
%\newmdtheoremenv[style=leftbar]{lemma}[theorem]{Lemma}
%\newmdtheoremenv[style=leftbar]{corollary}[theorem]{corollary}

\newtheorem{theorem}{Theorem}[section]
\newtheorem{proposition}[theorem]{Proposition}
\newtheorem{lemma}[theorem]{Lemma}
\newtheorem{corollary}[theorem]{Corollary}
\newtheorem{axiom}[theorem]{Axiom}
\newtheorem*{axiom*}{Axiom}

\surroundwithmdframed[style=leftbar]{theorem}
\surroundwithmdframed[style=leftbar]{proposition}
\surroundwithmdframed[style=leftbar]{lemma}
\surroundwithmdframed[style=leftbar]{corollary}
\surroundwithmdframed[style=leftbar]{axiom}
\surroundwithmdframed[style=leftbar]{axiom*}

\theoremstyle{definition}

\newtheorem*{definition}{Definition}
\surroundwithmdframed[style=leftbar]{definition}

\newtheorem*{slogan}{Slogan}
\newtheorem*{eg}{Example}
\newtheorem*{ex}{Exercise}
\newtheorem*{remark}{Remark}
\newtheorem*{notation}{Notation}
\newtheorem*{convention}{Convention}
\newtheorem*{assumption}{Assumption}
\newtheorem*{question}{Question}
\newtheorem*{answer}{Answer}
\newtheorem*{note}{Note}
\newtheorem*{application}{Application}

%operator macros

%basic
\DeclareMathOperator{\lcm}{lcm}

%matrix
\DeclareMathOperator{\tr}{tr}
\DeclareMathOperator{\Tr}{Tr}
\DeclareMathOperator{\adj}{adj}

%algebra
\DeclareMathOperator{\Hom}{Hom}
\DeclareMathOperator{\End}{End}
\DeclareMathOperator{\id}{id}
\DeclareMathOperator{\im}{im}
\DeclarePairedDelimiter{\generation}{\langle}{\rangle}

%groups
\DeclareMathOperator{\sym}{Sym}
\DeclareMathOperator{\sgn}{sgn}
\DeclareMathOperator{\inn}{Inn}
\DeclareMathOperator{\aut}{Aut}
\DeclareMathOperator{\GL}{GL}
\DeclareMathOperator{\SL}{SL}
\DeclareMathOperator{\PGL}{PGL}
\DeclareMathOperator{\PSL}{PSL}
\DeclareMathOperator{\SU}{SU}
\DeclareMathOperator{\UU}{U}
\DeclareMathOperator{\SO}{SO}
\DeclareMathOperator{\OO}{O}
\DeclareMathOperator{\PSU}{PSU}

%hyperbolic
\DeclareMathOperator{\sech}{sech}

%field, galois heory
\DeclareMathOperator{\ch}{ch}
\DeclareMathOperator{\gal}{Gal}
\DeclareMathOperator{\emb}{Emb}



%ceiling and floor
%https://tex.stackexchange.com/a/118217/26707
\DeclarePairedDelimiter\ceil{\lceil}{\rceil}
\DeclarePairedDelimiter\floor{\lfloor}{\rfloor}


\DeclarePairedDelimiter{\innerproduct}{\langle}{\rangle}

%\DeclarePairedDelimiterX{\norm}[1]{\lVert}{\rVert}{#1}
\DeclarePairedDelimiter{\norm}{\lVert}{\rVert}



%Dirac notation
%TODO: rewrite for variable number of arguments
\DeclarePairedDelimiterX{\braket}[2]{\langle}{\rangle}{#1 \delimsize\vert #2}
\DeclarePairedDelimiterX{\braketthree}[3]{\langle}{\rangle}{#1 \delimsize\vert #2 \delimsize\vert #3}

\DeclarePairedDelimiter{\bra}{\langle}{\rvert}
\DeclarePairedDelimiter{\ket}{\lvert}{\rangle}




%macros

%general

%divide, not divide
\newcommand*{\divides}{\mid}
\newcommand*{\ndivides}{\nmid}
%vector, i.e. mathbf
%https://tex.stackexchange.com/a/45746/26707
\newcommand*{\V}[1]{{\ensuremath{\symbf{#1}}}}
%closure
\newcommand*{\cl}[1]{\overline{#1}}
%conjugate
\newcommand*{\conj}[1]{\overline{#1}}
%set complement
\newcommand*{\stcomp}[1]{\overline{#1}}
\newcommand*{\compose}{\circ}
\newcommand*{\nto}{\nrightarrow}
\newcommand*{\p}{\partial}
%embed
\newcommand*{\embed}{\hookrightarrow}
%surjection
\newcommand*{\surj}{\twoheadrightarrow}
%power set
\newcommand*{\powerset}{\mathcal{P}}

%matrix
\newcommand*{\matrixring}{\mathcal{M}}

%groups
\newcommand*{\normal}{\trianglelefteq}
%rings
\newcommand*{\ideal}{\trianglelefteq}

%fields
\renewcommand*{\C}{{\mathbb{C}}}
\newcommand*{\R}{{\mathbb{R}}}
\newcommand*{\Q}{{\mathbb{Q}}}
\newcommand*{\Z}{{\mathbb{Z}}}
\newcommand*{\N}{{\mathbb{N}}}
\newcommand*{\F}{{\mathbb{F}}}
%not really but I think this belongs here
\newcommand*{\A}{{\mathbb{A}}}

%asymptotic
\newcommand*{\bigO}{O}
\newcommand*{\smallo}{o}

%probability
\newcommand*{\prob}{\mathbb{P}}
\newcommand*{\E}{\mathbb{E}}

%vector calculus
\newcommand*{\gradient}{\V \nabla}
\newcommand*{\divergence}{\gradient \cdot}
\newcommand*{\curl}{\gradient \cdot}

%logic
\newcommand*{\yields}{\vdash}
\newcommand*{\nyields}{\nvdash}

%differential geometry
\renewcommand*{\H}{\mathbb{H}}
\newcommand*{\transversal}{\pitchfork}
\renewcommand{\d}{\mathrm{d}} % exterior derivative

%number theory
\newcommand*{\legendre}[2]{\genfrac{(}{)}{}{}{#1}{#2}}%Legendre symbol


\graphicspath{{figures/}}

\renewcommand*{\P}{\mathbb{P}}
\DeclareMathOperator{\Pic}{Pic}

\begin{document}

\begin{titlepage}
  \begin{center}
    \includegraphics[width=0.6\textwidth]{logo.jpg}\par
    \vspace{1cm}
    {\scshape\huge Mathematics Tripos \par}
    \vspace{2cm}
    {\huge Part \npart \par}
    \vspace{0.6cm}
    {\Huge \bfseries \ntitle \par}
    \vspace{1.2cm}
    {\Large\nterm, \nyear \par}
    \vspace{2cm}
    
    {\large \emph{Lectures by } \par}
    \vspace{0.2cm}
    {\Large \scshape \nlecturer}
    
    \vspace{0.5cm}
    {\large \emph{Notes by }\par}
    \vspace{0.2cm}
    {\Large \scshape \href{mailto:\nauthoremail}{\nauthor}}
 \end{center}
\end{titlepage}

\tableofcontents

\section{Toric variety}

Roughly speaking, algebraic geometry studies the duality between polynomial equations and varieties. Let \(S = \C[x_1, \dots, x_n]\). Let \(I \subseteq S\) be an ideal generated by \(f_1, \dots, f_s\). We have polynomials of different flavours:

\subparagraph{monomial ideal}

A monomial ideal is an ideal \(I \subseteq S\) generated by monomials. The corresponding vanishing loci \(V(I)\) is a union of coordinate subspaces of \(\C^n\).

\begin{eg}\leavevmode
  \begin{itemize}
  \item \(x_1 = 0\) in \(S\), \(x_1x_2 = 0\) in \(\C[x_1, x_2]\) which is the union of two axes in \(\C^2\) so reducible.
  \item \(V(x_1^2)\) is the axis with ``additional structure'' as \((x_1^2)\) is not reduced.
  \end{itemize}
\end{eg}

Monomial ideals are thus the simplest kind of varieties and can be studied in a completely combinatorial way.

\paragraph{trinomial ideals}

A trinomial is a polynomial with at most three nonzero distinct monomial terms. Let \(I \subseteq S\) be an ideal generated by trinomials.

\begin{lemma}
  Any affine variety is the vanishing locus of a trinomial ideal.
\end{lemma}

\begin{remark}
  To be more precise, \(S/I\) as a ring is isomorphic to \(\C[x_1, \dots, C_N]/\mathfrak a\) where \(\mathfrak a \subseteq \C[x_1, \dots, x_n]\) is generated by trinomials.
\end{remark}

In other words, there is no way to distinguish a trinomial ideal and a variety.

\begin{proof}[Sketch]
  There is a trick to prove the lemma. Let \(I \subseteq S\) be an ideal. Let \(a_1 x^{\alpha_1} + \dots + a_\ell x^{\alpha_\ell} = 0\) be an equation determined by \(I\). For each such equation, introduce \((\ell - 3)\) new variables \(\{z_i\}_{i = 1}^{\ell - 3}\). Replace the equation with
  \begin{align*}
    a_1 x^{\alpha_1} + a_2 x^{\alpha_2} + z_1
    &= -z_1 + z_2 + a_3 x^{\alpha_3} \\
    &= -z_2 + z_3 + a_4 x^{\alpha_4} \\
    &= \cdots \\
    &= -z_{\ell - 3} + a_{\ell - 1} x^{\alpha_{\ell - 1}} + a_\ell x^{\alpha_\ell}
  \end{align*}
  Repeat for other equations.
\end{proof}

\paragraph{binomial ideals}

Let \(I \subseteq \C[x_1, \dots, x_n] = S\) be an ideal generated by binomials. Assume that \(I\) is prime (so \(V(I)\) is irreducible and reduced, therefore integral). Let \(X = V(I)\), say \(\dim X = m\). Then we have
\begin{itemize}
\item \(X\) contains \((\C^*)^m\) as a Zariski dense open subset.
\item The multiplication map \((\C^*)^m \times (\C^*)^m \to (\C^*)^m\) extends to a map \(X \times (\C^*)^m \to X\).
\end{itemize}

\begin{eg}\leavevmode
  \begin{itemize}
  \item \(\C^n\) generated by \(I = (0)\). It contains \((\C^*)^n\) as a dense open subset and we have multiplication
    \begin{align*}
      \C^n \times (\C^*)^n &\to \C^n \\
      ((z_1, \dots, z_n), (t_1, \dots, t_n)) &\mapsto (z_1t_1, \dots, z_nt_n)
    \end{align*}
  \item \(V(x_1 - x_2) \subseteq \C^3\) is isomorphic to \(\C^2\). \(\C^2\) contains a copy of \((\C^*)^2\), given by \(\{(t, t, t'): t, t' \in \C^*\}\).
  \item \(V(x_1^2 - x_2) \subseteq \C^3\) contains a copy of \((\C^*)^2\), given by \(\{(t, t^2, t'): t, t' \in \C\}\).
  \end{itemize}
\end{eg}

Note that the key fact here is that a monomial is a homomorphism \((\C^*)^n \to \C^*\) so the ideal \(I\) encoded the intersection of translates of kernels of homomorphisms.

\begin{definition}[algebraic torus]\index{algebraic torus}
  An \emph{algebraic torus} (over \(\C\)) of rank \(1\) is the group variety \(\mathbb G_m = \C^*\). A rank \(n\) algebraic torus is \(\mathbb G_m^n = (\C^*)^n\).
\end{definition}

\begin{remark}
  \(\mathbb G_m^n\) is a group variety as it is the loci \(V(x_1 \cdots x_{n + 1} - 1)\) and the group operations are all morphisms of varieties.
\end{remark}

When we say \(X\) is an affine variety, we mean \(X = V(I)\) where \(I \subseteq \C[x_1, \dots, x_n]\) and so \(X \subseteq \A^n_\C\). Usually we are not so interested in the embedding and only interested in the algebraic object. The coordinate ring \(S/I\) can have many presentations as a quotient of a polynomial ring.

\begin{notation}
  We write \(\Spec (S/J)\) to mean \(V(I)\), and when we write \(\Spec(R)\) we mean \(R = S/I\) and \(\Spec(R)\) is the space \(V(I)\).
\end{notation}

\begin{definition}[toric variety]\index{toric variety}
  An \emph{toric variety} \(X\) is an irreducible variety over \(\C\) of dimension \(\ell\) containing \((\C^*)^\ell\) as a dense open subset, and such that \(m: (\C^*)^\ell \times (\C^*)^\ell = (\C^*)^\ell\) extends to \(X \times (\C^*)^\ell = X\) (as a group action).
\end{definition}

\begin{remark}
  Since \(X\) is irreducible, a nonempty open set is dense.
\end{remark}

\begin{remark}
  A variety is a finite type scheme over \(\C\).
\end{remark}

\begin{eg}\leavevmode
  \begin{itemize}
  \item \((\C^*)^\ell\) for \(\ell \geq 0\).
  \item \(\C^n\).
  \item \((\C^*)^\ell \times \C^\ell\).
  \item \(\P^\ell\): pick coordinates \(x_0, \dots, x_\ell\). Then \((\C^*)^\ell = \{(1, t_1, \dots, t_\ell): t_i \in \C^*\} \embed \P^\ell\). Check it is open and dense: \(\P^n \supseteq \C^\ell = \{(x_0, \dots, x_\ell): x_0 \neq 0\}\) as an open dense subset, which further contains \((\C^*)^\ell\). Note that the torus depends on choice of coordinates: for example take two distinct \(p_1, p_2 \in \P^1\), there exists a change of coordinates (in \(\PGL(2)\)) such that \(p_1 \mapsto 0, p_2 \mapsto \infty\), so \(\P^1 \setminus \{p_1, p_2\} \cong \P^1 \setminus \{0, \infty\} \cong \C^*\).
  \item \(\P^{\ell_1} \times \dots \times \P^{\ell_k}\).
  \end{itemize}
\end{eg}

\begin{remark}
  There exist infintely many toric varieties in every dimension \(\ell \geq 2\). There exists \(X\) toric such that \(X\) is singular, or \(X\) could be compact but not projective. \(X\) could have torsion in \(\Pic(X)\). All of these can be constructed using combinatorial methods.
\end{remark}

\subsection{Two lattices}

We define two lattices. The first one is cocharacter \(N = \Hom(\C^*, (\C^*)^\ell)\) (as algebraic groups), which is isomorphic as an abstract group to \(\Z^\ell\). This is the isomorphism
\[
  (\lambda_1, \dots, \lambda_\ell) \mapsto (t \mapsto (t^{\lambda_1}, \dots, t^{\lambda_\ell})).
\]

The second is character \(M = \Hom((\C^*)^\ell, \C^*)\) which is also isomorphic to \(\Z^\ell\), given by
\[
  (u_1, \dots, u_\ell) \mapsto ((t_1, \dots, t_\ell) \mapsto t_1^{u_1} \cdots t_\ell^{u_\ell}).
\]

Let \(X\) be a toric variety with toris \(\mathbb G_m^\ell\) and lattices \(N\) and \(M\). Given \(\lambda \in N\), we obtain a morphism \(\lambda: \C^* \to (\C^*)^\ell \subseteq X\). We ask two questions: does \(\lim_{t \to \infty} \lambda(t)\) exist? If so what is it?

\begin{eg}
  Take \(X = \P^2\) with \((\C^*)^2 = \{(1, t_1, t_2): t_i \in \C^\times\}\). Then \(N \cong \Z^2\). The limit always exists as \(X = \P^2\) is proper. Let \(\lambda = (\lambda_1, \lambda_2) \in \Z^2\). There are seven cases:
  \begin{itemize}
  \item If \(\lambda_1, \lambda_2 > 0\) then \(\lambda(t) = (1, t^{\lambda_1}, t^{\lambda_2})\). Then \(\lim_{t \to 0} \lim(t) = (1, 0, 0)\).
  \item If \(\lambda_1 = 0, \lambda_2 > 0\) then \(\lim_{t \to 0} (1, 1, t^{\lambda_2}) = (1, 1, 0)\).
  \item \(\lambda_1 > \lambda_2, \lambda_2 < 0\). Then
    \[
      \lim_{t \to 0} (1, t^{\lambda_1}, t^{\lambda_2})
      = \lim_{t \to 0} (t^{-\lambda_2}, t^{\lambda_1 - \lambda_2}, 1)
      = (0, 0, 1)
    \]
  \item \dots
  \end{itemize}
  \begin{figure}[ht]
    \centering
    \def\svgwidth{.3\columnwidth}
    \input{figures/fan_P2.pdf_tex}
    \caption{Fan of \(\P^2\)}
  \end{figure}
  This is a \emph{fan} \(N \otimes \R = N_\R \cong \R^2\).
\end{eg}

\begin{figure}[ht]
  \centering
  \begin{minipage}{0.5\textwidth}
    \centering
    \def\svgwidth{.8\columnwidth}
    \input{figures/fan_P1.pdf_tex}
  \end{minipage}%
  \begin{minipage}{0.5\textwidth}
    \centering
    \def\svgwidth{.8\columnwidth}
    \input{figures/fan_P1P1.pdf_tex}
  \end{minipage}
  \caption{Fan of \(\P^1\) and \(\P^1 \times \P^1\)}
\end{figure}

Every toric variety \(X\) determines a \emph{fan} \(\Sigma \subseteq N_\R\). A fan is a union of \emph{cones} in \(N_\R\), where a cone is a set of the form
\[
  \operatorname{Cone}(S) = \{v \in N_\R: \sum_{i = 1}^r \nu_i \lambda_i, \nu_i \geq 0\}
\]
where \(S \subseteq N\) finite.






\iffalse
Multi-index notation: if \(\alpha = (\alpha_1, \dots, \alpha_n) \in \N^n\) then
\[
  X^\alpha = x_1^{\alpha_1} \cdots x_n^{\alpha_n}.
\]
\fi

\printindex
\end{document}
